% ============================================================================
% Abstract  (front matter, input before the table of contents)
% Structure after Jacquet: stakes and question / method / findings with
% numbers / implications. Four paragraphs, self-sufficient. No em dashes.
% ============================================================================

\section*{Abstract}

Europe's next wave of offshore wind moves from the mature North Sea into the
enclosed Baltic Sea, where wind farms generate at the same time and push down
the price they all sell at. The share of the average electricity price that a
wind farm actually receives is called its capture value, and it falls as more
wind is built. This thesis measures capture value for every electricity market
around the two seas from 2017 to 2024 and projects it to 2030 and 2040, using
open data throughout. It asks three questions: how the two seas compare today
and how they will compare once the planned wind farms are built, what causes
the loss of value and how much of it can be avoided by building the
surrounding system differently, and whether a model trained on the past can
predict a future this different, or whether a model that rebuilds the market
from its parts is required.

Two models are built and compared with each other. The first is a statistical
model, trained on eight years of observed prices at 31 offshore wind farms.
The second is a power-system model, PyPSA-Eur, which reconstructs the market
from its parts, the power plants, the transmission grid and the demand in
every hour, under the official European expansion scenarios for 2030 and 2040.
Both models are checked against observed prices before either is trusted.

Today the two regions are equal in terms of capture value: wind farms in both capture roughly 88 per cent
of the average electricity price. By 2040, with the planned buildout in place
and today's grid, Baltic farms fall to 47 per cent while North Sea farms hold
57, so the Baltic loses almost half of its present value. About half of that
loss is caused by congestion in the grid, and it can be removed: with
transmission expanded, the Baltic recovers to 66 per cent, level with the
North Sea again. The two models prove equally accurate where they can be
checked, and their estimates of capture value agree to within half a
percentage point, but beyond the observed record the statistical model returns
almost the same answer whichever future it is given, so the long-term
projection must come from the power-system model. Even with an expanded grid
the Baltic captures a quarter less than it does today, so what the planned
wind farms will be worth depends on the grid built alongside them.